
\documentclass[10pt,article,twocolumn]{report}
\usepackage[T1]{fontenc}
\usepackage{float}
\usepackage[utf8]{inputenc}
\usepackage{ctex}
\usepackage[utf8]{inputenc}
\usepackage{amsmath}
\usepackage{amssymb}
\usepackage{xcolor}
\usepackage{graphicx}
\usepackage{float}
\usepackage[utf8]{inputenc}
\usepackage[T1]{fontenc}
\usepackage{ctex}
\usepackage{amsmath}
\usepackage{amssymb}
\usepackage{xcolor}
\usepackage{graphicx}
\usepackage{float}
\usepackage{lmodern}
\usepackage{mathrsfs}
\usepackage{lipsum}
\usepackage{tikz}
\usepackage{braket}
\usepackage{adjustbox}
\usepackage{footmisc}
\usepackage{graphicx}
\usepackage{tcolorbox}
\setlength{\parindent}{0pt}
\title{谐振电路}
\author{张世航}
\begin{document}
	\maketitle
    \section{$I-f$串联谐振曲线}
    \begin{figure}[H]
    	\centering
    	\includegraphics[width=0.7\linewidth]{1}
    	\caption{$I-f$串联谐振曲线}
    	\label{fig:1}
    \end{figure}
   \section{谐振频率}
   对于串联谐振电路,我们有\[Z=\sqrt{R^2+(\omega L-\frac{1}{\omega C})^2}\]
   于是回路中的电流为\[I=\frac{U}{Z}=\frac{U}{\sqrt{R^2+(\omega L -\frac{1}{\omega C})^2}}\]
   于是我们发现在\[\omega L=\frac{1}{\omega C}\]
   时电路达到谐振状态.我们可以解出谐振角频率为\[\omega=\omega_0=\frac{1}{\sqrt{LC}}\]
   于是谐振频率为
   \begin{align*}
   	f&=\frac{1}{2\pi}\omega\\
   	&=\frac{1}{2\pi\sqrt{LC}}\\
   	&=\frac{1}{2\times\sqrt{0.1H\times4.7\times10^{-8}F}}\\
   	&=2322.690947Hz
   \end{align*}
   由于我们只有有限多个数据点,因此我们需要利用拟合来得到函数,进而通过梯度下降的方法得到极值.在这里我们利用导数的零点来求解最大值.\\
   \begin{figure}[H]
   	\centering
   	\includegraphics[width=0.7\linewidth]{5}
   	\caption{拟合曲线}
   	\label{fig:4}
   \end{figure}
   
   \subsection{$R=200\Omega$}
   \begin{figure}[H]
   	\centering
   	\includegraphics[width=0.7\linewidth]{2}
   	\caption{$R=200\Omega$}
   	\label{fig:2}
   \end{figure}
   得到极大值点为\[f=2323.4775454620694Hz\]
   \subsection{$R=500\Omega$}
   \begin{figure}[H]
   	\centering
   	\includegraphics[width=0.7\linewidth]{3}
   	\caption{$R=500\Omega$}
   	\label{fig:3}
   	得到极大值点为\[f=2322.7463520988094Hz\]
   \end{figure}
   
   很显然这两个值与我们得到的理论值相差很小.\\
   \section{通频带宽度}
   \begin{figure}[H]
   	\centering
   	\includegraphics[width=0.7\linewidth]{6}
   	\caption{频带宽度}
   	\label{fig:6}
   \end{figure}
   
   $R=200\Omega$的频带宽度:$f_{200\Omega}$=433.41852706111786 Hz\\
   $R=500\Omega$的频带宽度:$f_{500\Omega}$=905.8907175671729 Hz
   \section{串联谐振电路的Q值}
   \subsection{$R=200\Omega$的情况}
   理论值
   \begin{align*}
   	Q&=\sqrt{\frac{L}{{(R+R_S+R_L)}^2C}}\\
   	&=\sqrt{\frac{0.1}{(25+50+500)^2\times4.7\times10^{-8}}}\\
   	&=2.53
   \end{align*}
   而\[Q=\frac{U_C}{U}=2.544\]
   我们还可以利用频带宽度计算
   \[Q=\frac{f_0}{\Delta f}=\frac{2323.4775Hz}{905.891Hz}=2.56\]
   \subsection{$R=500\Omega$}
   理论值
   \begin{align*}
   	Q&=\sqrt{\frac{L}{{(R+R_S+R_L)}^2C}}\\
   	&=\sqrt{\frac{0.1}{(25+50+200)^2\times4.7\times10^{-8}}}\\
   	&=5.30
   \end{align*}
   而\[Q=\frac{U_C}{U}=5.28\]
   我们还可以利用频带宽度计算
   \[Q=\frac{f_0}{\Delta f}=\frac{2322.746Hz}{433.4185Hz}=5.36\]
\end{document}